\documentclass[12pt]{article}
\usepackage[utf8]{inputenc}
\usepackage[T1]{fontenc}
\usepackage[margin=1in]{geometry}
\usepackage{setspace}
\usepackage{xcolor}
\usepackage[colorlinks=true, urlcolor=blue!60!black, linkcolor=blue!60!black]{hyperref}
\setlength{\parindent}{0em}
\setlength{\parskip}{0.45em}
\renewcommand{\baselinestretch}{1.0}
\begin{document}
\begin{center}
{\Large \textit{Tinker v. Des Moines Independent Community School District}, 393 U.S. 503 (1969)\par}
\vspace{0.3em}
{\normalsize Case Brief --- One-Page Edition\par}
\end{center}

\textbf{Memory Jogger.} The black-armband case: student speech at the schoolhouse gate (\href{https://supreme.justia.com/cases/federal/us/393/503/}{393 U.S. 503}).

\textbf{Facts.} Three Des Moines students, including Mary Beth and John Tinker, wore black armbands in December 1965 to protest the Vietnam War. Principals had preemptively banned armbands, though other political symbols were allowed. The three refused to remove them and were suspended. Only a handful of 18,000 students wore armbands, and apart from a few hostile remarks there were no threats, violence, or disruption of schoolwork.

\textbf{Procedural History.} The district court upheld the ban; the en banc Eighth Circuit divided equally.

\textbf{Votes.} Reversed, 7--2. Fortas wrote for the Court; Stewart and White concurred; Black and Harlan dissented.

\textbf{Holding.} Schools may not punish passive political expression absent facts reasonably forecasting material and substantial disruption or invasion of others' rights.

\textbf{Reasoning.} The armbands were silent, near-pure speech; officials acted on undifferentiated fear and singled out one viewpoint while allowing campaign buttons. Black warned that the Court handed schools to their ``loudest-mouthed, but maybe not their brightest, students.''

\textbf{Examples.} \emph{Same-Side:} a quiet political T-shirt. \emph{Opposite-Side:} a neutral ban after documented fights. \emph{Fence-Sitter:} an off-campus post mocking a teacher.

\textbf{Critique.} Critics warn that the forecast test invites a heckler's veto and that the rights-of-others prong is undefined.

\textbf{Key Quotation.} ``It can hardly be argued that either students or teachers shed their constitutional rights to freedom of speech or expression at the schoolhouse gate.'' --- Justice Fortas, Opinion of the Court, 393 U.S. at 506.

\end{document}
